\section{Symbol-wise Source Coding}
\paragraph{Setup}
We denote the set of all symbols (the \textcolor{Blue}{alphabet}) by
$\mathcal{X}$. Consider a discrete (memoryless) source $X$ with PMF $P_X$ over
$\mathcal{X}$.\\
{\textcolor{ForestGreen}{Lowercase English text}}:
$\mathcal{X}=\{a,\ldots,z\}$, with $P_X(e)$ being the highest.

We encode each string $x\in \mathcal{X}^{*}$ into a decodeable binary string
$C(x)$, of length $\ell(x)$. The goal is to minimize the
\textcolor{Blue}{expected codeword length} of the code $C$:
\[L(C)=\mathbb{E}_{X\sim P_X}\left[ \ell\left(X\right) \right]=\sum_{x\in\mathcal{X}} P_X(x)\cdot\ell\left(x\right),\]
\vspace{-4mm}
\subsection{Symbol-Wise Coding}
\textcolor{Blue}{Symbol-wise coding} encodes \textcolor{Bittersweet}{each
symbol} $x\in\mathcal{X}$ into a codeword $C(x)$, and
\textcolor{Bittersweet}{encodes a string} $x_1\cdots x_n$ by concatenating the
codewords $C(x_1)\cdots C(x_n)$.

\paragraph{Decodeability Conditions}
\textcolor{Blue}{Non-singular codes}: A code is non-singular if all codewords
are distinct.
\[\forall x,x'\in \mathcal{X}, x\neq x' \implies C(x)\neq C(x').\] 

\textcolor{Blue}{Uniquely decodable codes}: A code is uniquely decodable if any
concatenation of codewords can be uniquely parsed into codewords. That is, for
distinct $x_1\cdots x_n,x'_1\cdots x'_m\in \mathcal{X}^{*}$,
\[C(x_1)\cdots C(x_{n})\neq C(x_1')\cdots C(x_m').\]
Uniquely decodeability is \textcolor{Bittersweet}{non-trivial to verify}.

\textcolor{Blue}{Prefix-free codes}: A code is
prefix-free if no codeword is a prefix of another codeword.
\[C\text{ is prefix-free}\implies C\text{ is uniquely decodeable}.\]
Additionally, prefix-free codes can be decoded
\textcolor{Bittersweet}{instantaneously}, without needing to wait for the
entire string to be received.

\paragraph{Kraft's inequality}
{Kraft's inequality} states that for any prefix-free code $C$,
\[\sum_{x\in \mathcal{X}} 2^{-\ell(x)}\le 1.\]
Any prefix-free code can be \textcolor{Bittersweet}{represented as a
\textbf{binary tree}}, where each codeword corresponds to a leaf. The depth of
the leaf is the codeword length. It follows, by random traversal of the tree,
that the \textcolor{Bittersweet}{probability of reaching a leaf is
$2^{-\ell(x)}$}.

Conversely, for any set of codeword lengths $\left\{\ell(x)\right\}_{x\in
\mathcal{X}}$ satisfying Kraft's inequality, there exists a prefix-free code
with these codeword lengths.

\paragraph{Entropy lower bound}
Entropy provides a fundamental \textcolor{Bittersweet}{limit on the expected
codeword length} of any uniquely decodable code:
\[L(C)\ge H(X),\]
with equality iff ${P_X(x)}=2^{-\ell(x)}$ for all $x\in \mathcal{X}$.

Using the entropy lower bound, \textcolor{ForestGreen}{\textbf{Shannon-Fano
coding}} is a simple method to construct a \textcolor{Bittersweet}{prefix-free
code for a given distribution} $P_X$.

Ordering the symbols in decreasing probability, define:
\[C(x)=\text{binary expansion of }\sum_{x'<x} P_X(x'),\]
truncated to $\ell(x)=\lceil -\log_2{P_X(x)}\rceil$ bits.

A Shannon-Fano coding satisfies the \textcolor{Bittersweet}{expected codeword
length} of
\[H(X)\le L(C_{SF})\le H(X)+1.\] 

\textcolor{Blue}{Mismatched coding}: If the true distribution is $P_X$ but
\textcolor{Bittersweet}{lengths are chosen according to $Q_X$}, the
Shannon-Fano code has expected codeword length
\[L(C_Q)=H(P_X)+D(P_X\|Q_X).\]

\subsection{Huffman Coding}
Huffman coding constructs an \textcolor{Bittersweet}{optimal prefix-free code}
for a given distribution $P_X$. The \textbf{algorithm} is as follows:
\begin{itemize}
  \item Start with a forest of $|\mathcal{X}|$ trees, each consisting of a single node corresponding to a symbol $x\in \mathcal{X}$, with weight $P_X(x)$.
  \item While there are more than one tree in the forest:
    \begin{itemize}
      \item Select the two roots with the smallest weights, say $T_1$ and $T_2$.
      \item Merge $T_1$ and $T_2$ into a new tree $T$, where $T_1$ and $T_2$ are the left and right subtrees of $T$, respectively. The weight of $T$ is the sum of the weights of $T_1$ and $T_2$.
    \end{itemize}
  \item The remaining tree in the forest is the Huffman tree.
\end{itemize}
At each node, assign a binary digit to each branch. The code for a given symbol
is the \textcolor{Bittersweet}{sequence of binary digits along the path} from
the root to the leaf corresponding to that symbol.
\vspace{-4mm}
\begin{center}
  \includegraphics[width=0.23\textwidth]{huffman.png}
\end{center}
\vspace{-4mm}
\textcolor{Blue}{Optimality}: The expected codeword length of the Huffman code is
\[L(C_H)=\min_{C\text{ prefix-free}} L(C).\]

\subsection{Extensions of Symbol-wise Coding}
\textcolor{Blue}{Block coding}: Instead of encoding each symbol separately, we
can \textcolor{Bittersweet}{\textbf{encode blocks} of $n$ symbols} together. We
consider instead the \textcolor{Bittersweet}{alphabet $\mathcal{X}^n$}, and the
distribution $P_{X^n}$.

This allows us to achieve expected codeword lengths closer to the entropy, as
$n$ increases. $\forall\epsilon>0$, there exists a block code $C$ s.t.:
\[L(C)<H(X)+\epsilon.\] 

\textcolor{ForestGreen}{\textbf{Advantages} of block coding:}
\begin{itemize}
  \item Encode sources with memory, by treating blocks as symbols.
    \textcolor{ForestGreen}{P(apple) > P(ealpp)} despite containing the same characters.
  \item Achieve expected codeword length arbitrarily close to the entropy. Even with
    independent symbols,
    \[NH(X)\le L(C)\le NH(X)+1.\] 
    \[H(X)\le \text{Average length per symbol}\le H(X)+\frac{1}{N}.\] 
\end{itemize}
\textcolor{Bittersweet}{\textbf{Disadvantages} of block coding:}
\begin{itemize}
  \item Determining the product distribution $P_{X^n}$ can be difficult,
    especially for sources with memory.
  \item Encoding and decoding complexity grows exponentially with $n$. Requires
    sorting $\left| \mathcal{X} \right| ^{n}$ symbols by probability, and
    storing the codebook of size $\left| \mathcal{X} \right| ^{n}$.
\end{itemize}


